\chapter{ROS2-Grundlagen}

\begin{lernziele}
Nach diesem Kapitel kennen Sie Nodes, Topics, Services und Actions und können diese Begriffe auf das Robotersystem anwenden.
\end{lernziele}

\section{Was ist ROS2?}
\gls{ros2} ist ein Framework für Robotiksoftware. Es stellt Kommunikationsmechanismen, Werkzeuge und Konventionen bereit. ROS2 ist kein Betriebssystem im klassischen Sinn, sondern eine Middleware und Entwicklungsumgebung.

\section{Nodes}
Ein Node ist ein ausführbares Programm mit einer klaren Aufgabe. Beispiel: Ein \texttt{camera\_node} liest Kamerabilder und veröffentlicht sie.

\section{Topics}
Topics sind Nachrichtenkanäle. Ein Node veröffentlicht Nachrichten, andere Nodes abonnieren sie. Diese Publish-Subscribe-Kommunikation passt sehr gut zu Sensordaten.

\section{Services}
Services sind Anfrage-Antwort-Kommunikation. Sie eignen sich für kurze, gezielte Abfragen, zum Beispiel \enquote{Setze Parameter}.

\section{Actions}
Actions sind für länger laufende Aufgaben gedacht. Navigation zu einem Ziel ist ein typischer Action-Anwendungsfall, weil der Vorgang länger dauert und Feedback liefern kann.

\section{Beispiel}
\begin{verbatim}
/camera/image_raw       Kamerabild
/scan                   Laserscan
/cmd_vel                Bewegungsbefehl
/battery_state          Akkustatus
/detected_objects       erkannte Objekte
\end{verbatim}

\begin{uebungbox}
Ordnen Sie den Funktionen Kamera, Navigation, Batterieüberwachung und Objekterkennung passende ROS2-Nodes und Topics zu.
\end{uebungbox}
